% closing_remarks.tex - Final chapter closing and resources
% ============================================================================

\section{Where to Go From Here}
\label{sec:next-steps}

This book gave you the essentials. The real learning happens when you start using Python for your own research. Here are resources to help you continue.

\subsection{Packages to Explore}

These libraries extend your toolkit for specific tasks. With what you've learned, you can teach yourself any of them from their documentation:

\begin{itemize}
    \item \textbf{seaborn} --- Statistical visualization built on Matplotlib. Less code, beautiful plots.
    \item \textbf{scikit-learn} --- Machine learning.
    \item \textbf{statsmodels} --- Statistical modeling: linear models, time series, hypothesis tests.
    \item \textbf{sympy} --- Symbolic mathematics: algebra, calculus, equation solving (Mathematica like).
    \item \textbf{lmfit} --- Curve fitting with parameter constraints, confidence intervals, \& model evaluation.
    \item \textbf{uncertainties} --- Automatic error propagation through calculations.
    \item \textbf{pint} --- Physical units handling (never confuse meters and feet again).
\end{itemize}

\subsection{Domain-Specific Libraries}

Search for Python packages in your field---they almost certainly exist:
\begin{itemize}
    \item \textbf{Chemistry:} RDKit, Open Babel, cclib, PyMOL
    \item \textbf{Biology:} Biopython, scikit-bio, scanpy
    \item \textbf{Physics/Astronomy:} Astropy, QuTiP, PyMC
    \item \textbf{Geoscience:} xarray, cartopy, MetPy
    \item \textbf{Image Analysis:} scikit-image, OpenCV, Pillow
\end{itemize}

\subsection{Improve Your Workflow}

As your projects grow, these practices will save you time and headaches:
\begin{itemize}
    \item Learn Git for version control
    \item Write tests for your code
    \item Create reusable packages from your scripts
    \item Explore Jupyter notebooks for interactive analysis
\end{itemize}

\subsection{Practice}

The best way to learn is to use Python for something you actually care about. Pick a dataset from your research. Automate something tedious. Recreate a figure from a paper. Break things, fix them, and learn by doing.

You now have enough foundation to teach yourself the rest.

\subsection{AI in the Scientific Workflow}
\label{sec:ai-usage}

\begin{note}
    Large Language Models (LLMs) are powerful assistive technologies, but they reward prepared minds disproportionately. This section offers suggestions for integrating them into your learning without undermining the skill development that makes you an independent researcher.
\end{note}

\noindent As you progress, you will inevitably encounter LLMs (AI chatbots) in your coding workflow. When used developmentally, they can accelerate your transition from syntax memorization to architectural thinking. Used prematurely, they can create an illusion of competence that fragments under scrutiny. The goal is to harness the former while avoiding the latter.

\subsubsection*{Understanding the Tool}
Think of an LLM as an \textit{ontological cartographer}---a system that maps relationships between concepts based on statistical patterns across text, rather than experiential understanding. 

You understand a ``reaction mechanism'' through laboratory experience, which carries a lot of information we never care to encode in language. The model understands ``reaction mechanism'' only in so far as can be derived from what we cared to capture using language. This is not a limitation to resent but a translation layer to manage. The model translates between linguistic representations of concepts; you must provide the physical-world verification.

\begin{ponder}
Think about how much is left unsaid in literature methods sections, or lab manuals. How could you expect a thing which lives in a world defined entirely by these incomplete captures to match your depth of comprehension of the topics?
\end{ponder}

\begin{warning}
An LLM cannot improve the quality of your question---it can only elaborate upon it. If you formulate a concept's question representation vaguely, that vagueness is a part of the concept you are invoking. So, when the LLM translates the concept to it's answer representation, quality predictably suffers. To win, you must understand the domain well enough to formulate good questions and evaluate the answers. For automated script writing, AI is most helpful when you are already proficient, not when you are a complete novice.
\end{warning}

\subsubsection*{A Developmental Framework}

Rather than viewing AI as forbidden or universally permitted, consider this progression:

\begin{enumerate}
    \item \textbf{Foundation Phase (Avoid generation, embrace explanation)} \\
    When learning a new library (e.g., your first RDKit script or lmfit model), do not ask the AI to write the code. Instead, write your own attempt, then use AI to:
    \begin{itemize}
        \item Explain error tracebacks in plain language
        \item Compare your approach to idiomatic patterns
        \item Suggest documentation sections you may have missed
    \end{itemize}
    
    \item \textbf{Apprenticeship Phase (Pair programming)} \\
    Once you can write working but inelegant code, use AI as a collaborative reviewer:
    \begin{itemize}
        \item Refactor working scripts for readability
        \item Vectorize slow loops you have already prototyped
        \item Generate unit tests for functions you have written
    \end{itemize}
    
    \item \textbf{Independence Phase (Managed delegation)} \\
    When you have the skill to write the code yourself but choose to delegate to save time, you are essentially trading the effort of writing for the effort of auditing. Because your prompts are condensed summaries of intent rather than literal instructions, the LLM must fill the "ambiguity gap"---the space between human language and executable code---with educated guesses. Delegation only works if you are skilled enough to recognise when those guesses miss the mark; if your words were precise enough to never be misunderstood, they would already be code. Coding without oversight is dangerous and deeply incompetent behaviour in professional environments.
    \begin{itemize}
        \item Generate boilerplate for familiar patterns, then audit line-by-line
        \item Prototype alternative algorithms you already understand conceptually
        \item Document code you have written (reverse the typical workflow)
    \end{itemize}
\end{enumerate}

\begin{tip}
\textbf{The explain-it-back Rule} \\
Before submitting any AI-assisted code, apply the laboratory notebook standard: could you explain this snippet to a labmate during group meeting? If not, treat it as a learning opportunity, not a submission. Open a separate file, break the code into chunks, and annotate each line until the logic is transparent.
\end{tip}

\subsubsection*{Mentorship and Cognitive Load}

Novice code written by humans has predictable error patterns---syntax mistakes, logical gaps, misconception signatures. These are diagnostically valuable; a mentor can identify the underlying confusion quickly, turning debugging into targeted teaching.

AI-generated novice code lacks these tell-tale patterns. It appears plausible but fails in subtle ways (e.g., using deprecated matplotlib arguments that still run but produce wrong chemical visualisations). This creates asymmetric debugging: you learn nothing, while your mentor expends disproportionate effort uncovering the error.

The solution is transparency, not avoidance. When seeking help with AI-assisted code:
\begin{itemize}
    \item Present your original attempt alongside the AI version
    \item Highlight specific lines you do not fully understand
    \item Ask: ``What conceptual gaps would cause this type of error?'' rather than just ``Fix this.''
\end{itemize}
This transforms the interaction from resentment-generating noise into a diagnostic dialogue about your reasoning.

\subsubsection*{Speed vs.\ Velocity}

It is tempting to view AI as a productivity multiplier, but in skill acquisition, velocity matters more than speed. Velocity implies direction; speed merely measures motion. 

Writing a script manually might take four hours; with AI, forty minutes. But if the manual process taught you to recognise array broadcasting errors, you have gained velocity---the next similar task takes thirty minutes regardless of AI assistance. If you skipped the struggle, you remain dependent on the tool for identical tasks forever.

\textit{Slow is smooth, and smooth is fast.} Master the concept at a pace that allows comprehension, \textit{then} use AI to accelerate implementation of what you genuinely understand. While AI can provide speed but YOU provide direction. Without your intentional direction guiding things, with AI you are simply directionless faster.

\subsubsection*{A Case Study in Confident Failure of LLMs}

While writing this book, as an exercise an LLM (Claude 4.5 Haiku, and Claude 4.6 Opus) were asked to produce performance comparisons like those in Section~\ref{sec:Numba}: four implementations of a simplified bimolecular reaction simulation (pure Python, unvectorized NumPy, vectorized NumPy, and Numba). The model produced plausible code and confidently declared vectorized NumPy the fastest. The initial timings appeared to confirm this. What followed was an instructive series of failures---not one mistake, but a chain of them, each requiring domain expertise to diagnose and resolve.

\medskip
\noindent\textbf{Act 1: The Hidden Bug.} The vectorized implementation computed pairwise distances between all molecules, but failed to exclude the matrix diagonal---the distance from each molecule to \emph{itself}, which is always zero. Every molecule ``reacted with itself'' on the first timestep, marking all 500 molecules as reacted instantly. The function then skipped the expensive distance computation for the remaining 9,999 timesteps, producing the illusion of speed (finishing in less than 0.01 seconds). None of this was obvious at a glance from the code alone; and the output \emph{looked} approximately correct.

It took careful proofing to catch the discrepancy: why did vectorized NumPy always report exactly 500 molecules reacted, while the other methods varied around 475$\pm$10 of 500? The model, when confronted with this inconsistency, initially blamed different random number generators, then suggested other irrelevant changes. Only after explicit, repeated, and pointed prompting about the distance matrix diagonal did it identify that issue. A practitioner with domain experience would likely not have made this error, and would easily have caught this during code review---excluding self-interactions in pairwise computations is not exotic; it is a basic requirement.

\medskip
\noindent\textbf{Act 2: The Overcorrection.} With the diagonal fixed (\py{np.fill\_diagonal(distances, np.inf)}), vectorized NumPy became the \emph{slowest} approach: 68 seconds vs.\ 1.6 for Numba. The model had constructed an entire pedagogical narrative around this finding, and recommended sweeping changes the book. However, the reason was straightforward: the vectorized version allocated a full 500$\times$500 distance matrix at every time-step, while the loop-based versions skipped already-reacted molecules with \py{continue}, so their workload shrank over time. The model missed this entirely, because understanding \emph{why} an algorithm is slow requires reasoning about memory allocation patterns and workload dynamics---not just pattern-matching against common advice. If a LLM considers these factors or not in an answer is really a matter of luck. Without specific prompting to bring these factors into it's attention is fairly unlikely to be consistently addressed.

\medskip
\noindent\textbf{Act 3: The Fix the Model Couldn't Find.} The solution was to subset to only active (unreacted) molecules before computing the distance matrix, so that the array shrank as molecules reacted. This brought vectorized NumPy to $\sim$1 second---still the fastest single-invocation approach, but at a more reasonable margin. All fixes and optimisations had to be suggested by the human, otherwise would be missed by models. It requires understanding of \emph{both} the algorithmic structure (pairwise computations on a shrinking set) \emph{and} NumPy's performance characteristics (small array operations are fast; huge allocations are expensive). The model's repeated failure to suggest this---despite it being a natural and well-known pattern---illustrates how brittle AI can be even with simply domain shifts. There are other domains where this kind of thing is far more commonplace and so I doubt a LLM coding assistant would ever make a mistake such as this.

\medskip
Several things are worth noting:

\begin{enumerate}
    \item The cascade was insidious: a subtle bug produced plausible timings, which were used to construct a compelling but false narrative. Each layer of interpretation moved further from the underlying error, making it harder to trace back. At every stage, the model's output \emph{read} convincingly.
    \item Even after the bug was found, getting the right answer required a second insight---the active subsetting optimisation---that the model did not volunteer. Fixing one error does not guarantee correctness; it often reveals the next layer of complexity.
    \item The bug was not exotic. Excluding self-interactions and subsetting active particles are standard considerations in any pairwise simulation. The model lacked the physical intuition to catch the first and the performance intuition to suggest the second.
    \item At every stage, it took human domain expertise to drive improvement and correction. The model provides syntax; the human provides understanding.
	\item Using AI it took close to two hours and all of a days worth of a session limit's worth of tokens to reproduce what took 20 minutes to write otherwise. This would likely have been faster had full scope of problem been supplied in the initial query (making sure to ignore diagonal for reactions, managing memory and throughput by masking matrix)
\end{enumerate}

This illustrates a fundamental asymmetry. LLMs perform best on well-documented, widely-discussed problems. As you approach the specifics of your work---novel analyses, unusual data structures, domain-specific edge cases---the model's training data thins out, and the probability of confident but incorrect output increases. This is precisely where your expertise matters most, and precisely where uncritical delegation is most dangerous.

\begin{ponder}
The closer your work is to the cutting edge of your field, the less likely it is that a sufficiently capable model exists to handle it reliably. At that point, the only thing that can catch the errors is your own domain knowledge. How do you develop that knowledge if you delegate the work that would build it?
\end{ponder}

\begin{warning}
\textbf{Reputation and Ownership} \\
In professional contexts---industry R\&D, graduate research groups, collaborative publications---your code represents your scientific rigour. Submitting opaque AI-generated scripts that fail under review signals unreliable judgement. Conversely, demonstrating fluency in explaining, verifying, and modifying AI-assisted workflows signals mature computational literacy. The tool is neutral; your stewardship of it is not.
\end{warning}

\subsection{Where to Get Help}
If you do get stuck (and you will be---everyone does) in order of quality:
\begin{itemize}
    \item \textbf{Official Documentation} --- NumPy, Pandas, and Matplotlib all have excellent tutorials and example galleries.
	\item \textbf{Your colleagues}---They will often know more about the specific problems you're trying to solve than most others.
    \item \textbf{Stack Overflow} --- Search before asking. Your error message has almost certainly been answered.
    \item \textbf{GitHub Issues} --- For package-specific bugs or feature questions.
    \item \textbf{Reddit} --- r/learnpython is friendly to beginners; r/Python for general discussion.
	\item \textbf{Chat-bots}---Undeniably useful, even if these can be problematic. Use with caution, and don't accept code you don't understand. These are a double-edged blade, and you WILL get cut using these.
\end{itemize}

\vfill
\begin{center}
    \textit{Happy Coding!}
\end{center}